\chapter{绪论}\label{chap:introduction}

异常检测（Anomaly Detection, AD）作为数据挖掘与机器学习领域的核心研究任务，旨在从海量高维数据中识别与正常分布存在显著偏离的少数偏离样本。异常检测在金融欺诈识别、医疗辅助诊断、IT系统网络入侵检测以及工业传感器故障诊断等领域具有重大的科学研究价值与广泛的工业落地前景。然而，在学术界与工业界现有的研究体系中，异常检测正面临着前所未有的“跨模态割裂”与“数据标注污染”双重严峻挑战。

首先，不同模态（表格、时序、图、图像、文本）的数据物理分布与关联特征存在巨大的几何差异，导致异常检测在各模态下的最优算法体系长期处于分裂状态。其次，在真实的工业落地场景中，海量数据的精准标注极其昂贵、繁琐且极易出错，使得异常检测的训练集往往混入大量的噪声标签（如对称标签翻转、未标注的隐蔽异常）。这种严重的“数据/标签污染”会使经典的强判别模型产生严重的过拟合。

针对上述学术与工业界的共性痛点，本文构建了一个统一的\textbf{跨模态抗污染异常检测大一统基准平台}，跨越表格（Tabular）、时序（TimeSeries）、图（Graph）、图像（CV）与文本（NLP）五大模态，通过部署 14 种经典与前沿算法在 29 个真实数据集上进行大规模抗噪防御消融大评测，并针对 20\% 的极端对称标签翻转提出了一种极具工程可行性的主动防御包装器（Robust Defense Wrapper），系统探究了抗噪防御策略的性能边界。

\section{研究背景与意义}

\subsection{跨模态异常检测的科学价值}
异常检测算法的最优选择具有极强的数据驱动和模态依赖性。传统的异常检测研究长期处于“各自为战”的烟囱式发展状态。例如，在表格（Tabular）数据中，由于样本间不具备明显的空间或时间关联，基于树集成（如 LightGBM \cite{ke2017lightgbm}、XGBoost \cite{chen2016xgboost}）和邻近度量（如 KNN \cite{zhao2019pyod}）的算法往往表现出最优的几何分类面表达能力；在时序（TimeSeries）数据中，异常通常隐含在时间戳之间的长短期时间依赖或局部突发波动之中，因此基于距离（如 MatrixProfile）和时序预测重构（如 LSTM-AE）的无监督时序技术成为主流；而在图（Graph）数据中，网络拓扑结构和节点特征的协同融合（如 GCN \cite{gadbench2024}）是捕获结构性异常的关键。

建立一个大一统的跨模态异常检测评估大底座，不仅能够系统评估不同算法族群（树模型、经典无监督模型、前沿大模型与深度循环模型）在异构数据模态上的底层泛化能力，更能够为学术界打破模态割裂、探索具有统一架构的跨模态异常检测大模型（Foundation Models for Anomaly Detection）提供坚实的理论支撑。

\subsection{数据与标签污染的工程隐忧}
在实际工程落地中，数据污染和标签噪声（Label Noise）几乎不可避免。由于人工多轮核验的局限性、系统噪声或黑客的恶意投毒，数据集中经常会引入比例不一的标签噪声，其中最经典的数学建模为对称标签翻转（Symmetric Label Flips）\cite{frenay2013labelnoise}。
对称标签翻转会导致以下两个致命问题：
\begin{enumerate}
    \item \textbf{过度拟合噪声特征}：像神经网络（MLP）这样的强判别分类器，在训练过程中会极力拟合被错误标记为异常的正常样本（或者被错误标记为正常的异常样本），导致决策边界发生严重的非线性扭曲，在干净测试集上的性能发生断崖式下跌。
    \item \textbf{决策边界塌陷}：部分简单线性模型或弱分类器在面对混淆的噪声标签时，由于不具备非线性的高阶特征过滤能力，其在损失函数反向传播时的梯度更新会相互抵消，导致最终分类器偏向于预测先验众数，使异常检测决策面发生根本性崩溃。
\end{enumerate}
因此，如何设计一种无需重新训练底层分类器、具备即插即用特性（Plug-and-Play）的主动鲁棒防御包装器（Robust Defense Wrapper）\cite{jiang2020trim}，成为工业界异常检测系统实现长效高鲁棒运行的生命线。

\section{本文研究内容与主要学术贡献}

本文旨在打破现存的多模态壁垒，针对标签污染场景提出一种主动鲁棒防御的学术基准。我们的工作主要从以下四个大规模阶梯式实验展开：

\begin{enumerate}
    \item \textbf{第一组实验（无噪基准评估, Exp1）}：建立 5 大模态、29 个数据集大一统的无噪基准，对比有监督与无监督在干净状态下的上限。
    \item \textbf{第二组实验（噪声退化评估, Exp2）}：在 0\% 至 20\% 的标签污染梯度下，描绘各算法性能衰退的渐进性退化曲线，评估固有鲁棒性。
    \item \textbf{第三组实验（模态难度分析, Exp3）}：跨越模态障碍，量化不同模态的分类面搜索难度，定位各算法的跨模态全局排名。
    \item \textbf{第四组实验（主动防御消融大评测, Exp4）}：设计并部署一个主动式防御包装器，横向对比 9 种无监督去噪器（IQR, IForest, AutoEncoder 等）在 Trim（剔除）和 Flip（翻转）策略下的消融表现，并科学定义防御算法的失效边界。
\end{enumerate}

\subsection{本文的核心学术贡献}
本文的核心学术贡献主要体现在以下三个方面：
\begin{enumerate}
    \item \textbf{大一统跨模态数据流基准平台}：首次实现了一个可以同时加载、适配并以统一标准化接口输出表格、时序、图、CV 和 NLP 数据包的集成流水线，极大降低了异常检测的多模态迁移壁垒。
    \item \textbf{提出高效的主动鲁棒防御包装器并验证其优越性}：证实了基于样本剪裁策略的无监督去噪防御包装器（特别是 \texttt{IQR\_Trim} 策略）在 20\% 的极端标签噪声下，不仅几乎不损害干净数据（0\% 污染）时的表现，更能够跨算法大幅拉升异常检测的 AUC-ROC（使 MLP 性能暴涨 \textbf{2.52\%}），展现出其强大的工程鲁棒拯救能力。
    \item \textbf{定义了鲁棒防御 wrapper 的“失效边界”}：首次在学术上深度剖析并证实了基于样本剪裁的防御策略在预训练元学习大模型（如 TabPFN \cite{hollmann2022tabpfn}）和线性分类模型（如 Logistic Regression）下的副作用及内在数学原理，为工业界防噪决策树、深度循环网络等非线性模型的部署提供了精准而理性的“失效边界（No Free Lunch）”理论边界。
\end{enumerate}

\subsection{论文的组织结构}
本学位论文（报告）的后续章节组织结构安排如下：
\begin{itemize}
    \item 第二章（相关工作）对异常检测算法体系、标签噪声学习理论和现有基准研究进行系统的文献综述，确立本文的研究坐标系。
    \item 第三章（系统方法与架构）详细阐述本文提出的大一统数据管道架构、对称标签污染模型的数学刻画以及基于无监督样本去噪的主动防御包装器技术。
    \item 第四章（实验设计与实施）介绍涉及的 29 个真实数据集属性、 baseline 超参配置、软硬件开发平台和复现细节。
    \item 第五章（实验结果与分析）系统陈述 Exp1 至 Exp4 四组大型实验的实测成绩，插入动态包络图、消融柱状图进行多维度横向剖析。
    \item 第六章（深度讨论与局限性剖析）针对 LSTM-Sup 的类别不平衡崩溃以及 TabPFN、LR 的防御失效边界开展严谨的底层学术探讨与根源剖析。
    \item 第七章（结论与展望）总结本文的全部研究成果，指出非对称标签噪声防御和图对比防御等未来的探索方向。
\end{itemize}


